% GENERATED FILE. Edit canonical lessons, then run npm run generate:books.
% canonical-source-hash: fnv1a64:9e090a7733906b48
% canonical-lessons: MR-C181-sun, MR-C181-navardev, MR-C181-navri, MR-C181-vidhva, MR-C181-yas

\chapter{Daughter-in-law, Bridegroom, Bride, Widow, Success}
\label{ch:mr-daughter-in-law-bridegroom-bride-widow-success}
\hlchaptermodality{181}

\begin{chapteropening}
\textbf{By the end of this chapter:} \emph{I can say a daughter-in-law, a bridegroom, a bride, a widow and success.}

Complete the last lesson of chapter 181: I can say a daughter-in-law, a bridegroom, a bride, a widow and success.
\end{chapteropening}

\section[\texorpdfstring{\mr{सून} (sūn)}{सून (sūn)}]{\mr{सून} (sūn) — a daughter-in-law}

\label{lesson:MR-C181-sun}



\begin{quote}
Before the new one: say the Marathi for a journalist, then the Marathi for an architect.
\end{quote}

\subsection*{You'll want to know: \mr{सून}}
\textbf{\mr{सून}} — \emph{sūn} — \textquotedblleft{}a daughter-in-law\textquotedblright{}.

\begin{grammarlens}[title={What you've built}]
One more word for this chapter.
\end{grammarlens}

\subsection*{Your turn}
\begin{itemize}
\raggedright
  \item \emph{Say it:} \emph{sūn}
  \item \emph{Say it:} \emph{sūn}, once more
  \item \emph{Say it:} say \emph{sūn}
  \item \emph{Recall:} say the Marathi for a journalist, then the Marathi for an architect, then say \emph{sūn} again
\end{itemize}

\begin{tcolorbox}[breakable,colback=teal!4,colframe=teal!35!black,title={Before you move on}]
What does \mr{सून} mean? (\textquotedblleft{}A daughter-in-law\textquotedblright{}.) Say it once more.
\end{tcolorbox}
\section[\texorpdfstring{\mr{नवरदेव} (navardev)}{नवरदेव (navardev)}]{\mr{नवरदेव} (navardev) — a bridegroom}

\label{lesson:MR-C181-navardev}



\begin{quote}
Before the new one: say the Marathi for an architect, then the Marathi for a daughter-in-law.
\end{quote}

\subsection*{You'll want to know: \mr{नवरदेव}}
\textbf{\mr{नवरदेव}} — \emph{navardev} — \textquotedblleft{}a bridegroom\textquotedblright{}.

\begin{grammarlens}[title={What you've built}]
One more word for this chapter.
\end{grammarlens}

\subsection*{Your turn}
\begin{itemize}
\raggedright
  \item \emph{Say it:} \emph{navardev}
  \item \emph{Say it:} \emph{navardev}, once more
  \item \emph{Say it:} say \emph{navardev}
  \item \emph{Recall:} say the Marathi for an architect, then the Marathi for a daughter-in-law, then say \emph{navardev} again
\end{itemize}

\begin{tcolorbox}[breakable,colback=teal!4,colframe=teal!35!black,title={Before you move on}]
What does \mr{नवरदेव} mean? (\textquotedblleft{}A bridegroom\textquotedblright{}.) Say it once more.
\end{tcolorbox}
\section[\texorpdfstring{\mr{नवरी} (navrī)}{नवरी (navrī)}]{\mr{नवरी} (navrī) — a bride}

\label{lesson:MR-C181-navri}



\begin{quote}
Before the new one: say the Marathi for a daughter-in-law, then the Marathi for a bridegroom.
\end{quote}

\subsection*{You'll want to know: \mr{नवरी}}
\textbf{\mr{नवरी}} — \emph{navrī} — \textquotedblleft{}a bride\textquotedblright{}.

\begin{grammarlens}[title={What you've built}]
One more word for this chapter.
\end{grammarlens}

\subsection*{Your turn}
\begin{itemize}
\raggedright
  \item \emph{Say it:} \emph{navrī}
  \item \emph{Say it:} \emph{navrī}, once more
  \item \emph{Say it:} say \emph{navrī}
  \item \emph{Recall:} say the Marathi for a daughter-in-law, then the Marathi for a bridegroom, then say \emph{navrī} again
\end{itemize}

\begin{tcolorbox}[breakable,colback=teal!4,colframe=teal!35!black,title={Before you move on}]
What does \mr{नवरी} mean? (\textquotedblleft{}A bride\textquotedblright{}.) Say it once more.
\end{tcolorbox}
\section[\texorpdfstring{\mr{विधवा} (vidhvā)}{विधवा (vidhvā)}]{\mr{विधवा} (vidhvā) — a widow}

\label{lesson:MR-C181-vidhva}



\begin{quote}
Before the new one: say the Marathi for a bridegroom, then the Marathi for a bride.
\end{quote}

\subsection*{You'll want to know: \mr{विधवा}}
\textbf{\mr{विधवा}} — \emph{vidhvā} — \textquotedblleft{}a widow\textquotedblright{}.

\begin{grammarlens}[title={What you've built}]
One more word for this chapter.
\end{grammarlens}

\subsection*{Your turn}
\begin{itemize}
\raggedright
  \item \emph{Say it:} \emph{vidhvā}
  \item \emph{Say it:} \emph{vidhvā}, once more
  \item \emph{Say it:} say \emph{vidhvā}
  \item \emph{Recall:} say the Marathi for a bridegroom, then the Marathi for a bride, then say \emph{vidhvā} again
\end{itemize}

\begin{tcolorbox}[breakable,colback=teal!4,colframe=teal!35!black,title={Before you move on}]
What does \mr{विधवा} mean? (\textquotedblleft{}A widow\textquotedblright{}.) Say it once more.
\end{tcolorbox}
\section[\texorpdfstring{\mr{यश} (yaś)}{यश (yaś)}]{\mr{यश} (yaś) — success}

\label{lesson:MR-C181-yas}



\begin{quote}
Before the new one: say the Marathi for a bride, then the Marathi for a widow.
\end{quote}

\subsection*{You'll want to know: \mr{यश}}
\textbf{\mr{यश}} — \emph{yaś} — \textquotedblleft{}success\textquotedblright{}.

\begin{grammarlens}[title={What you've built}]
One more word for this chapter.
\end{grammarlens}

\subsection*{Your turn}
\begin{itemize}
\raggedright
  \item \emph{Say it:} \emph{yaś}
  \item \emph{Say it:} \emph{yaś}, once more
  \item \emph{Say it:} say \emph{yaś}
  \item \emph{Recall:} say the Marathi for a bride, then the Marathi for a widow, then say \emph{yaś} again
\end{itemize}

\begin{tcolorbox}[breakable,colback=teal!4,colframe=teal!35!black,title={Before you move on}]
What does \mr{यश} mean? (\textquotedblleft{}Success\textquotedblright{}.) Say it once more.
\end{tcolorbox}
